% --- CORRECCIÓN IMPORTANTE EN LA PRIMERA LÍNEA ---
% Agregamos "xcolor=table" para que funcione \rowcolor sin errores
\documentclass[aspectratio=169, 10pt, xcolor=table]{beamer}

% --- Configuración del Tema (estilo profesional como el ejemplo) ---
\usetheme{Madrid}
\usecolortheme{dolphin}

% --- Paquetes necesarios ---
\usepackage[utf8]{inputenc}
\usepackage[spanish]{babel}
\usepackage{amsmath, amssymb}
\usepackage{booktabs}
\usepackage{graphicx}
\usepackage{listings}
\usepackage{tikz}
\usetikzlibrary{arrows.meta, positioning, shapes.geometric}

% --- Ruta de Imágenes (crea la carpeta "imagenes" y coloca ahí tus archivos) ---
\graphicspath{{./imagenes/}}

% --- Estilo para código Python ---
\definecolor{codegreen}{rgb}{0,0.6,0}
\definecolor{codegray}{rgb}{0.5,0.5,0.5}
\definecolor{codepurple}{rgb}{0.58,0,0.82}
\definecolor{backcolour}{rgb}{0.95,0.95,0.92}
\lstset{
    language=Python,
    backgroundcolor=\color{backcolour},
    commentstyle=\color{codegreen},
    keywordstyle=\color{blue},
    numberstyle=\tiny\color{codegray},
    stringstyle=\color{codepurple},
    basicstyle=\ttfamily\scriptsize,
    breaklines=true,
    frame=single,
    showstringspaces=false
}

% --- Pie de página personalizado ---
\makeatletter

\setbeamertemplate{footline}{
  \leavevmode%
  \hbox{%
  \begin{beamercolorbox}[wd=.333333\paperwidth,ht=2.25ex,dp=1ex,center]{author in head/foot}%
    \usebeamerfont{author in head/foot}\insertshortauthor
  \end{beamercolorbox}%
  \begin{beamercolorbox}[wd=.333333\paperwidth,ht=2.25ex,dp=1ex,center]{title in head/foot}%
    \usebeamerfont{title in head/foot}Sesión 5
  \end{beamercolorbox}%
  \begin{beamercolorbox}[wd=.333333\paperwidth,ht=2.25ex,dp=1ex,right]{date in head/foot}%
    \usebeamerfont{date in head/foot}Machine Learning \hfill \insertframenumber/\inserttotalframenumber \hspace*{0.3cm}
  \end{beamercolorbox}}%
  \vskip0pt%
}

\makeatother

% --- Datos de la portada ---
\title[SESIÓN 5]{\textbf{Sesión 5}}
\subtitle{Clasificación II: KNN, Naive Bayes y SVM}
\author[Prof. C. Sánchez]{Carlos César Sánchez Coronel}
\institute{\textbf{MACHINE LEARNING}}
\date{}

% Logo opcional (descomentar si tienes la imagen)
% \titlegraphic{\includegraphics[height=1.5cm]{logo.png}}

% --- Eliminar la barra de navegación (opcional) ---
\setbeamertemplate{navigation symbols}{}

% --- Índice al inicio de cada sección ---
\AtBeginSection[]{
  \begin{frame}
    \frametitle{Contenido}
    \tableofcontents[currentsection]
  \end{frame}
}

% ============================================================================
% INICIO DEL DOCUMENTO
% ============================================================================
\begin{document}

% --- Portada ---
\begin{frame}
    \titlepage
\end{frame}

% --- Índice general ---
\begin{frame}{Contenido}
    \tableofcontents
\end{frame}

% ============================================================================
% SECCIÓN 1: INTRODUCCIÓN
% ============================================================================
\section{Introducción}

\begin{frame}{Más allá de la regresión logística}
    \begin{itemize}
        \item La regresión logística es un modelo lineal y probabilístico.
        \item Muchos problemas requieren enfoques diferentes:
        \begin{itemize}
            \item Relaciones no lineales.
            \item Datos con estructuras específicas (texto, imágenes).
            \item Necesidad de interpretabilidad probabilística.
        \end{itemize}
        \item En esta sesión veremos tres algoritmos clásicos:
        \begin{enumerate}
            \item \textbf{KNN}: basado en similitud geométrica.
            \item \textbf{Naive Bayes}: basado en probabilidad condicional.
            \item \textbf{SVM}: basado en optimización de márgenes.
        \end{enumerate}
    \end{itemize}
\end{frame}

% ============================================================================
% SECCIÓN 2: K-VECINOS MÁS CERCANOS (KNN)
% ============================================================================
\section{K-Vecinos más cercanos (KNN)}

\begin{frame}{KNN: Fundamento geométrico}
    \begin{block}{Idea principal}
        Puntos cercanos en el espacio de características tienden a tener la misma clase.
    \end{block}
    \begin{itemize}
        \item No paramétrico, basado en instancias (``perezoso'').
        \item Para un nuevo punto $\mathbf{x}_q$:
        \begin{enumerate}
            \item Buscar los $k$ vecinos más cercanos en el conjunto de entrenamiento.
            \item Clasificación: voto mayoritario.
            \item Regresión: promedio de los valores.
        \end{enumerate}
    \end{itemize}
    % Basado en Pregunta 1
    % IMAGEN: Ejemplo de clasificación KNN en 2D
\end{frame}

\begin{frame}{Métricas de distancia}
    \begin{columns}[t]
        \column{0.5\textwidth}
        \begin{block}{Euclidiana}
            \[
            d(\mathbf{x},\mathbf{y}) = \sqrt{\sum_{j=1}^p (x_j - y_j)^2}
            \]
        \end{block}
        \begin{block}{Manhattan}
            \[
            d(\mathbf{x},\mathbf{y}) = \sum_{j=1}^p |x_j - y_j|
            \]
        \end{block}
        \column{0.5\textwidth}
        \begin{block}{Minkowski (general)}
            \[
            d(\mathbf{x},\mathbf{y}) = \left(\sum_{j=1}^p |x_j - y_j|^r\right)^{1/r}
            \]
        \end{block}
        \begin{block}{Coseno}
            \[
            d = 1 - \frac{\mathbf{x}\cdot\mathbf{y}}{\|\mathbf{x}\|\|\mathbf{y}\|}
            \]
        \end{block}
    \end{columns}
    \vspace{0.2cm}
    \begin{itemize}
        \item Manhattan es más robusta a outliers que Euclidiana.
        \item Coseno es útil para datos de alta dimensionalidad (texto).
    \end{itemize}
    % Basado en Pregunta 2
\end{frame}

\begin{frame}{Influencia del parámetro $k$}
    \begin{columns}[t]
        \column{0.5\textwidth}
        \begin{block}{$k$ pequeño (ej. $k=1$)}
            \begin{itemize}
                \item Modelo muy flexible.
                \item Alta varianza (sensible al ruido).
                \item Riesgo de sobreajuste.
            \end{itemize}
        \end{block}
        \column{0.5\textwidth}
        \begin{block}{$k$ grande}
            \begin{itemize}
                \item Modelo más estable.
                \item Aumenta el sesgo.
                \item Puede suavizar en exceso los límites.
            \end{itemize}
        \end{block}
    \end{columns}
    \vspace{0.3cm}
    \begin{exampleblock}{Selección de $k$}
        Validación cruzada para elegir el $k$ que minimice el error.
    \end{exampleblock}
    % Basado en Pregunta 1
    % IMAGEN: Gráfico de error vs k
\end{frame}

\begin{frame}{Importancia del escalado}
    \begin{alertblock}{¡Crítico en KNN!}
        KNN se basa en distancias. Si las variables tienen diferentes escalas, las de mayor magnitud dominarán el cálculo.
    \end{alertblock}
    \begin{columns}[t]
        \column{0.5\textwidth}
        \textbf{Estandarización (Z-score)}
        \[
        x' = \frac{x - \mu}{\sigma}
        \]
        \column{0.5\textwidth}
        \textbf{Normalización (Min-Max)}
        \[
        x' = \frac{x - x_{\min}}{x_{\max} - x_{\min}}
        \]
    \end{columns}
    \vspace{0.2cm}
    \begin{itemize}
        \item Se prefiere estandarización cuando hay outliers.
        \item En scikit-learn: \texttt{StandardScaler} o \texttt{MinMaxScaler}.
    \end{itemize}
    % Basado en Pregunta 2
\end{frame}

\begin{frame}{Maldición de la dimensionalidad}
    \begin{block}{Problema}
        Al aumentar el número de dimensiones $p$, los datos se vuelven dispersos y las distancias entre puntos tienden a ser similares.
    \end{block}
    \begin{itemize}
        \item KNN pierde capacidad para encontrar vecinos realmente cercanos.
        \item Soluciones:
        \begin{itemize}
            \item Reducción de dimensionalidad (PCA, selección de características).
            \item Métricas de distancia alternativas (coseno).
            \item Métodos de aproximación (LSH).
        \end{itemize}
    \end{itemize}
    % Basado en Pregunta 3
    % IMAGEN: Ilustración de la maldición
\end{frame}

\begin{frame}{Caso de uso: Sistemas de recomendación}
    \begin{exampleblock}{Recomendación basada en usuarios}
        \begin{enumerate}
            \item Representar usuarios por vectores de valoraciones.
            \item Para un usuario activo, encontrar los $k$ usuarios más similares (distancia de coseno o correlación).
            \item Recomendar ítems que esos vecinos hayan valorado positivamente y el activo no haya visto.
        \end{enumerate}
    \end{exampleblock}
    \begin{itemize}
        \item Problemas de escalabilidad con millones de usuarios $\rightarrow$ usar índices aproximados (LSH) o clustering previo.
        \item Cold start: para nuevos usuarios, usar popularidad o contenido.
    \end{itemize}
    % Basado en Pregunta 13
\end{frame}

\begin{frame}[fragile]{Implementación en scikit-learn}
    \begin{lstlisting}
from sklearn.neighbors import KNeighborsClassifier

knn = KNeighborsClassifier(n_neighbors=5, metric='euclidean')
knn.fit(X_train, y_train)
y_pred = knn.predict(X_test)
    \end{lstlisting}
    \begin{itemize}
        \item Parámetros importantes: \texttt{n\_neighbors}, \texttt{metric}, \texttt{weights} (unifom o distance).
        \item Para regresión: \texttt{KNeighborsRegressor}.
    \end{itemize}
\end{frame}

% ============================================================================
% SECCIÓN 3: NAIVE BAYES
% ============================================================================
\section{Naive Bayes}

\begin{frame}{Teorema de Bayes}
    \begin{block}{Fórmula}
        \[
        P(C_k|\mathbf{x}) = \frac{P(C_k) P(\mathbf{x}|C_k)}{P(\mathbf{x})}
        \]
        donde:
        \begin{itemize}
            \item $P(C_k)$: probabilidad a priori de la clase $k$.
            \item $P(\mathbf{x}|C_k)$: verosimilitud (probabilidad de los datos dada la clase).
            \item $P(C_k|\mathbf{x})$: probabilidad a posteriori.
        \end{itemize}
    \end{block}
    \begin{itemize}
        \item Clasificamos eligiendo la clase con mayor probabilidad a posteriori.
        \item El denominador $P(\mathbf{x})$ es constante y puede ignorarse.
    \end{itemize}
    % Basado en Pregunta 4
\end{frame}

\begin{frame}{Supuesto de independencia condicional}
    \begin{alertblock}{``Naive'' (ingenuo)}
        Se asume que las características son condicionalmente independientes dada la clase:
        \[
        P(\mathbf{x}|C_k) = \prod_{j=1}^p P(x_j|C_k)
        \]
    \end{alertblock}
    \begin{itemize}
        \item Este supuesto rara vez se cumple en la realidad.
        \item Sin embargo, simplifica enormemente la estimación y el modelo funciona sorprendentemente bien, especialmente en texto.
    \end{itemize}
    % Basado en Pregunta 4
\end{frame}

\begin{frame}{Variantes de Naive Bayes}
    \begin{columns}[t]
        \column{0.33\textwidth}
        \begin{block}{GaussianNB}
            Para características \textbf{continuas}:
            \[
            P(x_j|C_k) \sim \mathcal{N}(\mu_{kj}, \sigma_{kj}^2)
            \]
        \end{block}
        \column{0.33\textwidth}
        \begin{block}{MultinomialNB}
            Para \textbf{frecuencias} (conteo de palabras):
            \[
            P(\mathbf{x}|C_k) \propto \prod_j \theta_{kj}^{x_j}
            \]
        \end{block}
        \column{0.33\textwidth}
        \begin{block}{BernoulliNB}
            Para \textbf{binarias} (presencia/ausencia):
            \[
            P(x_j|C_k) = \theta_{kj}^{x_j}(1-\theta_{kj})^{1-x_j}
            \]
        \end{block}
    \end{columns}
    % Basado en Pregunta 5
\end{frame}

\begin{frame}{Suavizado Laplace}
    \begin{problem}
        Si una palabra nunca aparece en la clase ``spam'' en entrenamiento, su probabilidad condicional sería cero, anulando el producto.
    \end{problem}
    \begin{block}{Solución: Suavizado Laplace (add-1 smoothing)}
        \[
        P(x_j|C_k) = \frac{\text{freq}(x_j, C_k) + \alpha}{\sum_{j'} \text{freq}(x_{j'}, C_k) + \alpha |V|}
        \]
        con $\alpha = 1$ típicamente.
    \end{block}
    \begin{itemize}
        \item Evita probabilidades cero.
        \item Asigna una pequeña probabilidad a palabras no vistas.
    \end{itemize}
    % Basado en Pregunta 6
\end{frame}

\begin{frame}{Caso de uso: Clasificación de textos (spam)}
    \begin{exampleblock}{Filtro de spam}
        \begin{itemize}
            \item Características: presencia/frecuencia de palabras.
            \item Se usa \texttt{MultinomialNB} o \texttt{BernoulliNB}.
            \item Independencia condicional: asume que la probabilidad de ``viagra'' es independiente de ``oferta'' dado que es spam (falso, pero funciona).
            \item Rápido, escalable a miles de palabras.
        \end{itemize}
    \end{exampleblock}
    % Basado en Pregunta 4 y 12
\end{frame}

\begin{frame}[fragile]{Implementación en scikit-learn}
    \begin{lstlisting}
from sklearn.naive_bayes import GaussianNB, MultinomialNB, BernoulliNB

# Para características continuas
gnb = GaussianNB()
gnb.fit(X_train, y_train)

# Para frecuencias (texto)
mnb = MultinomialNB(alpha=1.0)
mnb.fit(X_train_counts, y_train)

# Para binarias
bnb = BernoulliNB(alpha=1.0)
bnb.fit(X_train_binary, y_train)
    \end{lstlisting}
\end{frame}

% ============================================================================
% SECCIÓN 4: MÁQUINAS DE VECTORES DE SOPORTE (SVM)
% ============================================================================
\section{Máquinas de Vectores de Soporte (SVM)}

\begin{frame}{SVM: Margen máximo}
    \begin{block}{Idea principal}
        Encontrar el hiperplano que separa las clases con el mayor margen posible.
    \end{block}
    \begin{itemize}
        \item Hiperplano: $\mathbf{w}^T\mathbf{x} + b = 0$.
        \item Decisión: $\text{sign}(\mathbf{w}^T\mathbf{x} + b)$.
        \item Margen: distancia mínima del hiperplano a los puntos más cercanos = $2/\|\mathbf{w}\|$.
        \item Maximizar el margen $\equiv$ minimizar $\frac{1}{2}\|\mathbf{w}\|^2$.
    \end{itemize}
    % Basado en Pregunta 7
    % IMAGEN: SVM con margen
\end{frame}

\begin{frame}{SVM linealmente separable}
    \begin{block}{Problema de optimización primal}
        \[
        \min_{\mathbf{w},b} \frac{1}{2}\|\mathbf{w}\|^2 \quad \text{sujeto a} \quad y_i(\mathbf{w}^T\mathbf{x}_i + b) \ge 1 \ \forall i
        \]
    \end{block}
    \begin{itemize}
        \item Los puntos con $y_i(\mathbf{w}^T\mathbf{x}_i + b) = 1$ son los \textbf{vectores soporte}.
        \item Sólo ellos determinan el hiperplano; el resto pueden eliminarse.
    \end{itemize}
    % Basado en Pregunta 7
\end{frame}

\begin{frame}{Margen suave y parámetro $C$}
    Para datos no separables linealmente, se introducen variables de holgura $\xi_i$:
    \[
    \min_{\mathbf{w},b,\boldsymbol{\xi}} \frac{1}{2}\|\mathbf{w}\|^2 + C\sum_{i=1}^n \xi_i
    \]
    sujeto a $y_i(\mathbf{w}^T\mathbf{x}_i + b) \ge 1 - \xi_i$, $\xi_i \ge 0$.
    \begin{itemize}
        \item $C$: controla la penalización por errores.
        \item $C$ grande $\rightarrow$ margen estrecho, pocos errores (riesgo de sobreajuste).
        \item $C$ pequeño $\rightarrow$ margen amplio, más errores permitidos (modelo simple).
    \end{itemize}
    % Basado en Pregunta 8
    % IMAGEN: Efecto de C
\end{frame}

\begin{frame}{Truco del kernel}
    \begin{block}{Idea}
        Proyectar los datos a un espacio de mayor dimensión donde sean separables linealmente, sin calcular explícitamente la transformación $\phi(\mathbf{x})$.
    \end{block}
    \begin{itemize}
        \item Se usa una función kernel $K(\mathbf{x}_i, \mathbf{x}_j) = \phi(\mathbf{x}_i)^T \phi(\mathbf{x}_j)$.
    \end{itemize}
    \begin{columns}[t]
        \column{0.33\textwidth}
        \textbf{Lineal}
        \[
        K(\mathbf{x}_i,\mathbf{x}_j) = \mathbf{x}_i^T \mathbf{x}_j
        \]
        \column{0.33\textwidth}
        \textbf{Polinomial}
        \[
        K(\mathbf{x}_i,\mathbf{x}_j) = (\gamma \mathbf{x}_i^T \mathbf{x}_j + r)^d
        \]
        \column{0.33\textwidth}
        \textbf{RBF}
        \[
        K(\mathbf{x}_i,\mathbf{x}_j) = \exp(-\gamma \|\mathbf{x}_i - \mathbf{x}_j\|^2)
        \]
    \end{columns}
    % Basado en Pregunta 9
\end{frame}

\begin{frame}{Parámetros clave en SVM con kernel RBF}
    \begin{columns}[t]
        \column{0.5\textwidth}
        \begin{block}{$\gamma$ (gamma)}
            Controla la influencia de un solo punto.
            \begin{itemize}
                \item $\gamma$ grande: solo puntos cercanos influyen (sobreajuste).
                \item $\gamma$ pequeño: modelo más suave (subajuste).
            \end{itemize}
        \end{block}
        \column{0.5\textwidth}
        \begin{block}{$C$}
            Penalización por error.
            \begin{itemize}
                \item Interacción con $\gamma$: ambos controlan la complejidad.
            \end{itemize}
        \end{block}
    \end{columns}
    \vspace{0.3cm}
    \begin{exampleblock}{Búsqueda de hiperparámetros}
        GridSearchCV o RandomSearchCV con validación cruzada.
    \end{exampleblock}
    % Basado en Pregunta 10
\end{frame}

\begin{frame}{Vectores soporte}
    \begin{itemize}
        \item Son los puntos que quedan sobre el margen o del lado incorrecto.
        \item Definen el hiperplano; si se eliminan, el modelo cambia.
        \item En problemas reales, pueden interpretarse como casos ``difíciles'' o representativos.
    \end{itemize}
    % Basado en Pregunta 14
    % IMAGEN: Vectores soporte destacados
\end{frame}

\begin{frame}{Caso de uso: Reconocimiento facial}
    \begin{exampleblock}{Clasificación de caras}
        \begin{itemize}
            \item Imágenes convertidas en vectores de píxeles (alta dimensionalidad).
            \item SVM con kernel RBF puede encontrar límites no lineales.
            \item Los vectores soporte son las imágenes más representativas o atípicas.
        \end{itemize}
    \end{exampleblock}
    % Basado en Pregunta 9
\end{frame}

\begin{frame}[fragile]{Implementación en scikit-learn}
    \begin{lstlisting}
from sklearn.svm import SVC

# SVM lineal
svm_lin = SVC(kernel='linear', C=1.0)
svm_lin.fit(X_train, y_train)

# SVM con kernel RBF
svm_rbf = SVC(kernel='rbf', C=1.0, gamma='scale')
svm_rbf.fit(X_train, y_train)
    \end{lstlisting}
    \begin{itemize}
        \item Para obtener probabilidades: \texttt{probability=True} (más lento).
    \end{itemize}
\end{frame}

% ============================================================================
% SECCIÓN 5: COMPARACIÓN DE ALGORITMOS
% ============================================================================
\section{Comparación}

\begin{frame}{Comparación: KNN vs Naive Bayes vs SVM}
    \begin{table}
        \centering
        \renewcommand{\arraystretch}{1.3}
        \begin{tabular}{l|c|c|c}
            \toprule
            \textbf{Característica} & \textbf{KNN} & \textbf{Naive Bayes} & \textbf{SVM} \\
            \midrule
            Interpretabilidad & Baja & Alta & Media (lineal) / Baja (RBF) \\
            Entrenamiento & Ninguno (perezoso) & Rápido & Medio (lineal) / Lento (RBF) \\
            Predicción & Lento & Rápido & Rápido \\
            Manejo de no linealidades & Sí (distancia) & No & Sí (con kernel) \\
            Escalado necesario & Sí & No (pero mejora) & Sí \\
            Alta dimensionalidad & Sufre & Robusto & Robusto (con kernel lineal) \\
            \bottomrule
        \end{tabular}
    \end{table}
    % Basado en Pregunta 11 y 17
\end{frame}

% ============================================================================
% SECCIÓN 6: CASO DE ESTUDIO INTEGRADO
% ============================================================================
\section{Caso de estudio: Clasificación de noticias}

\begin{frame}{Contexto y datos}
    \begin{block}{20 Newsgroups}
        \begin{itemize}
            \item $\approx$ 20,000 documentos en 20 categorías.
            \item Para simplificar, usamos dos categorías (ej. ``medicina'' vs ``espacio'').
        \end{itemize}
    \end{block}
    \begin{exampleblock}{Preprocesamiento}
        \begin{enumerate}
            \item Tokenización, eliminación de stopwords.
            \item Vectorización: CountVectorizer o TF-IDF.
            \item Reducción de vocabulario (palabras más frecuentes).
            \item División entrenamiento/prueba.
        \end{enumerate}
    \end{exampleblock}
    % Basado en Pregunta 12 y 19
\end{frame}

\begin{frame}{Modelado y comparación}
    Modelos a probar:
    \begin{enumerate}
        \item KNN con $k=5$, distancia euclidiana (o coseno).
        \item Multinomial Naive Bayes (con suavizado Laplace).
        \item SVM lineal.
        \item SVM con kernel RBF (con búsqueda de $C$ y $\gamma$).
    \end{enumerate}
    \vspace{0.3cm}
    \textbf{Resultados esperados:}
    \begin{itemize}
        \item MultinomialNB: muy bueno, rápido.
        \item SVM lineal: también excelente, especialmente con TF-IDF.
        \item SVM RBF: puede mejorar si la frontera no es lineal.
        \item KNN: puede sufrir por alta dimensionalidad y ser lento.
    \end{itemize}
\end{frame}

\begin{frame}{Análisis e interpretación}
    \begin{itemize}
        \item \textbf{Naive Bayes}: permite ver las palabras más probables por categoría (inspeccionar $\theta_{kj}$).
        \item \textbf{SVM lineal}: los coeficientes indican importancia de las palabras (para kernel lineal).
        \item \textbf{SVM con RBF}: menos interpretable; se pueden usar técnicas como SHAP o LIME.
        \item \textbf{KNN}: se pueden examinar los vecinos de un documento para entender la decisión.
    \end{itemize}
    % Basado en Pregunta 19
\end{frame}

% ============================================================================
% SECCIÓN 7: CASO AVANZADO: DIAGNÓSTICO MÉDICO (Pregunta 20)
% ============================================================================
\section{Caso avanzado: Diagnóstico médico}

\begin{frame}{Diseño del sistema}
    \begin{itemize}
        \item 100,000 pacientes, 50 variables (edad, presión, colesterol, ECG...).
        \item Objetivo: clasificar enfermedad cardíaca (binario).
        \item Requisitos: interpretable, alto recall (detectar enfermedad).
    \end{itemize}
    \begin{block}{Pipeline de preprocesamiento}
        \begin{enumerate}
            \item Imputar nulos (mediana).
            \item Detectar outliers (IQR) y decidir tratamiento.
            \item Estandarizar (Z-score).
            \item Selección de características (importancia de árboles, correlación).
        \end{enumerate}
    \end{block}
    % Basado en Pregunta 20
\end{frame}

\begin{frame}{Selección de modelo y ajuste}
    \begin{columns}[t]
        \column{0.5\textwidth}
        \textbf{Opciones:}
        \begin{itemize}
            \item KNN: $k$ óptimo por validación.
            \item GaussianNB: rápido, interpretable.
            \item SVM lineal: interpretable (coeficientes).
            \item SVM RBF: mejor rendimiento, menos interpretable.
        \end{itemize}
        \column{0.5\textwidth}
        \textbf{Para maximizar recall:}
        \begin{itemize}
            \item Ajustar umbral de decisión (ej. 0.3).
            \item Usar pesos de clase (class\_weight='balanced').
            \item Sobremuestreo (SMOTE) con cuidado.
        \end{itemize}
    \end{columns}
\end{frame}

\begin{frame}{Evaluación y monitoreo}
    \begin{itemize}
        \item \textbf{Métricas:} matriz de confusión, recall, precisión, F1, AUC-ROC, curva Precision-Recall.
        \item Priorizar recall, pero controlar precisión para evitar falsos positivos.
        \item \textbf{Interpretabilidad:} usar SHAP para explicar predicciones de SVM RBF.
        \item \textbf{Monitoreo en producción:} detectar data drift con PSI, KS, y reentrenar si es necesario.
    \end{itemize}
\end{frame}

% ============================================================================
% SECCIÓN 8: CONCLUSIÓN
% ============================================================================
\section{Conclusión}

\begin{frame}{Lo que hemos aprendido}
    \begin{exampleblock}{Conceptos clave}
        \begin{itemize}
            \item \textbf{KNN}: clasificador basado en instancias, sensible a escala y dimensionalidad.
            \item \textbf{Naive Bayes}: modelo probabilístico con supuesto de independencia, muy efectivo en texto.
            \item \textbf{SVM}: maximiza el margen, con truco del kernel para no linealidades.
            \item Comparación de fortalezas y debilidades según el problema.
            \item Aplicaciones a recomendación, spam, reconocimiento facial, diagnóstico médico.
        \end{itemize}
    \end{exampleblock}
    \begin{block}{Conexión con siguientes sesiones}
        Estos conceptos se extienden a métodos ensemble (Random Forest, Gradient Boosting) y redes neuronales.
    \end{block}
\end{frame}

% --- Diapositiva final ---
\begin{frame}
    \centering
    \Huge \textbf{¡Gracias!}

\end{frame}

\end{document}